\section*{Chapitre \ref{ch:m1:margin} --- Marge}

\begin{solution}{exo:m1:margin:1}
$-44.25$ points $\times\ \$50 \times 25 = -\$55\,312.50$~: le compte paie.
\end{solution}

\begin{solution}{exo:m1:margin:2}
Les fonds propres tombent à $120\,000 - 16\,000 = \$104\,000$, en dessous
de la \omterm{def:m1:margin:bond}{marge de maintenance} de $5 \times 21\,000 = \$105\,000$~: un \omterm{def:m1:financing:margincall}{appel de marge}
pour rétablir le niveau initial de $5 \times 23\,100 = \$115\,500$. Le client doit déposer 11\,500~\$.
\end{solution}

\begin{solution}{exo:m1:margin:3}
Un tiers de l'\omterm{def:m1:margin:scan}{intervalle de balayage} vaut 115 points, soit 5\,750~\$. Scénarios 1--2~: 0, 0. Puis
(perte positive) 3--4~: $-5\,750$~; 5--6~: $+5\,750$~; 7--8~: $-11\,500$~; 9--10~:
$+11\,500$~; 11--12~: $-17\,250$~; 13--14~: $+17\,250$~; 15~: $-15\,525$~; 16~:
$+15\,525$ ($3 \times 17\,250 \times 30\%$). La marge correspond à la plus grande perte~:
17\,250~\$, scénarios 13 et 14.
\end{solution}

\begin{solution}{exo:m1:margin:4}
Le portefeuille est vendeur d'options, donc vendeur de volatilité~: chaque scénario impair
ajoute quatre points de volatilité et fait monter la valeur de ce qu'il a vendu.
Le scénario 16 est pire que le 15, car le portefeuille est aussi acheteur de 2 contrats à terme~:
une baisse pénalise à la fois les contrats à terme et les puts vendus, tandis qu'en cas de hausse le
gain des contrats à terme compense une partie de la perte sur les calls vendus. Le scénario 13 (intervalle complet
à la baisse, volatilité à la hausse) serait contraignant si les options avaient un strike plus proche de la
monnaie, de sorte que l'essentiel de leur perte survienne dans un même \omterm{def:m1:margin:scan}{intervalle de balayage}, ou si les
scénarios extrêmes pesaient pour une fraction plus faible.
\end{solution}

\begin{solution}{exo:m1:margin:5}
$213\,250 - 107\,770 = \$105\,480$, soit 49~\%. À 5~\% par an~: 5\,274~\$ par an pour ce
petit portefeuille et, plus important encore, une taille de position deux fois plus petite pour le même
capital.
\end{solution}

\begin{solution}{exo:m1:margin:6}
$2.1\% \times \sqrt2 = 2.97\%$~; $2.1\% \times \sqrt5 = 4.70\%$. Plus élevé~: la volatilité
se regroupe en grappes et les rendements en période de crise sont positivement autocorrélés, et la
liquidation elle-même fait bouger les prix. Plus faible~: retour à la moyenne sur certains produits, et
un portefeuille partiellement couvert dès le premier jour ne porte pas tout son risque pendant
cinq jours.
\end{solution}

\begin{solution}{exo:m1:margin:7}
Hausses sur dix jours~: 109~\%, 203~\% et 94~\%. Moyennes en période calme~: 3.0~\% (avec filtre) et
4.5~\% (avec plancher à un poids de 25~\%). Un poids de 47~\% ramène la hausse sur dix jours
à 50~\%, pour une marge en période calme de 5.8~\% de la position, presque
le double du chiffre sans plancher~: la stabilité s'achète avec du collatéral qui reste
inactif la plupart du temps.
\end{solution}

\begin{solution}{exo:m1:margin:8}
En cas de \emph{hausse} de 10~\%, les contrats à terme perdent 20~millions de dollars, à régler en espèces dès le
lendemain matin, tandis que les 20~millions de dollars gagnés sur les actions sont latents~: le \omterm{def:m1:financing:pb}{prime
broker} prêtera sur ces gains, mais avec une décote et selon son propre calendrier, et pas du tout
s'il réduit son effet de levier au même moment. En cas de baisse, les espèces
arrivent de la chambre de compensation, mais le \omterm{def:m1:financing:pb}{prime broker} peut lancer un \omterm{def:m1:financing:margincall}{appel de marge} sur les
actions. Dans les deux sens, la \omterm{def:m1:clearing-and-settlement:margin}{marge initiale} sur les contrats à terme augmente avec
la volatilité. La neutralité de marché est une affirmation sur le P\&L~; la marge est une affirmation
sur les espèces, jambe par jambe, plateforme par plateforme.
\end{solution}

\begin{solution}{pb:m1:margin:1}
\textbf{1.} $400 \times 3\,300 \times 50 = \$66$ millions~; marge $4.5\% =
\$2.97$ millions.
\textbf{2.} $400 \times 50 \times 800 = \$16$ millions.
\textbf{3.} $400 \times 2\,500 \times 50 \times 9\% = \$4.5$ millions~: $+\$1.53$
millions, bien que le notionnel ait baissé d'un quart.
\textbf{4.} 17.53~millions de dollars, soit 26.6~\% du notionnel initial.
\textbf{5.} Il dispose de 12~millions de dollars. Il doit en réunir 5.5~millions en vendant
des obligations, au moment où elles sont les plus difficiles à vendre, ou réduire la position aux
plus bas.
\textbf{6.} Le contrat à terme tombe à 2\,488.75~: $400 \times 50 \times 261.25 =
\$5.225$ millions.
\textbf{7.} $400 \times 50 \times 165 = \$3.3$ millions, exigibles dans l'heure~:
espèces placées dans un fonds obligataire dont le règlement intervient en deux jours ne comptent pas. Un \omterm{def:m1:margin:intraday}{appel de marge
intrajournalier} manqué constitue un défaut comme les autres.
\textbf{8.} De $6\% \times 55$ millions $= \$3.3$ millions à $7.5\% \times
49.775$ millions $= \$3.733$ millions~: $+\$433\,125$.
\textbf{9.} $5\,225\,000 + 433\,125 = \$5\,658\,125$, dont 3.3~millions de dollars à
midi.
\textbf{10.} Le modèle filtré augmente le plus (la marge triple en dix jours)~; le
modèle à plancher augmente le moins en termes relatifs, le modèle historique simple le
moins en valeur absolue, car il a à peine bougé.
\textbf{11.} Les modèles filtré et à plancher, qui sont alors proches de
9~\% de la position~; la fenêtre historique simple reste proche de 5~\% alors que
la volatilité quotidienne est de 3.5~\%, si bien qu'un mouvement sur deux jours à 99~\% vaut environ 11~\%~: elle est
sous-margée de moitié.
\textbf{12.} Elle exige une marge supérieure de moitié en années calmes, quand les membres ne voient
aucun risque, les concurrents annoncent des marges plus faibles, et le coût est visible alors que
le bénéfice ne l'est pas.
\textbf{13.} La seconde, si le fonds dimensionne ses positions en conséquence~: ses \omterm{def:m1:financing:margincall}{appels de marge}
en pleine crise sont proportionnellement identiques, mais le fonds a toujours détenu moins
d'effet de levier et davantage de collatéral depuis le début. La première est moins chère chaque jour sauf
un.
\textbf{14.} Chaque position acheteuse et chaque position vendeuse déposent une \omterm{def:m1:clearing-and-settlement:margin}{marge initiale}~; doubler
le taux double le collatéral immobilisé à la chambre de compensation des deux côtés de
chaque contrat, alors que la \omterm{def:m1:clearing-and-settlement:margin}{marge de variation} déplace simplement des espèces des perdants vers
les gagnants.
\textbf{15.} Pas avec les espèces disponibles. Le gain sur les actions est chez un \omterm{def:m1:financing:pb}{prime broker}, qui
décide du montant qui peut être retiré~; la perte sur les contrats à terme est exigible en espèces à la
chambre de compensation le matin. La couverture supprime la perte, mais non l'\omterm{def:m1:financing:margincall}{appel de marge}.
\textbf{16.} Les espèces nécessaires pour tenir la position pendant un scénario de stress~: un mois
de \omterm{def:m1:clearing-and-settlement:margin}{marge de variation} défavorable à une volatilité stressée, auquel s'ajoute la hausse de
\omterm{def:m1:clearing-and-settlement:margin}{marge initiale} selon le modèle de la chambre de compensation dans le même scénario,
face aux espèces et aux actifs réellement éligibles et disponibles en quelques heures.
\textbf{17.} Les chambres de compensation acceptent une liste de collatéraux, principalement des espèces et
des obligations d'État, assorties de décotes~; les obligations d'entreprise sont souvent inéligibles ou
lourdement décotées, et la \omterm{def:m1:clearing-and-settlement:margin}{marge de variation} n'est constituée que d'espèces.
\textbf{18.} Les cours des obligations baissent, les fonds qui les détiennent font face à des rachats et à leurs
propres \omterm{def:m1:financing:margincall}{appels de marge}, la volatilité augmente, et les modèles de
la \cref{fig:m1:margin:procyclical} relèvent à nouveau les marges~: la boucle que le
terme \omterm{def:m1:margin:procyclical}{procyclicité} désigne.
\textbf{19.} \textbf{1.53~millions de dollars.}
\textbf{20.} La solvabilité porte sur la valeur des actifs face aux passifs~;
un \omterm{def:m1:financing:margincall}{appel de marge} porte sur des espèces dans la bonne devise, à la bonne banque, avant dix
heures.
\end{solution}

\begin{solution}{iq:m1:margin:1}
La \omterm{def:m1:clearing-and-settlement:margin}{marge de variation} règle en espèces le gain ou la perte du jour, de sorte qu'aucune dette
ne s'accumule entre le membre et la chambre de compensation~; elle passe des perdants aux
gagnants. La \omterm{def:m1:clearing-and-settlement:margin}{marge initiale} est un collatéral constitué contre la perte qui pourrait survenir
entre un défaut et la clôture des positions du défaillant~; elle est restituée lorsque
la position est clôturée. La première relève du P\&L, la seconde d'un dépôt dont la
taille peut être modifiée par la chambre de compensation.
\iqlookfor{le rôle de chacune, et le fait que la \omterm{def:m1:clearing-and-settlement:margin}{marge initiale} peut changer sans aucune modification
de la position.}
\end{solution}

\begin{solution}{iq:m1:margin:2}
Pour chaque produit, la chambre de compensation fixe un prix et un \omterm{def:m1:margin:scan}{intervalle de balayage} de volatilité.
Chaque contrat possède un tableau de risque~: sa perte selon seize scénarios, des variations de prix
égales aux tiers de l'intervalle, avec une volatilité à la hausse et à la baisse, auxquelles s'ajoutent
deux mouvements extrêmes pris en compte partiellement. On multiplie par les positions, on somme par scénario,
et l'on retient le pire~: c'est le \omterm{def:m1:margin:scan}{risque de balayage}. On ajoute une charge pour les spreads entre échéances,
on applique un minimum par option vendue, et l'on retranche les crédits entre produits corrélés.
\iqlookfor{les scénarios extrêmes et la raison du minimum sur les options
vendues.}
\end{solution}

\begin{solution}{iq:m1:margin:3}
Les marges qui augmentent en situation de stress ajoutent à la demande de liquidités au
moment où celles-ci sont les plus rares. Outils~: un plancher issu d'une période de stress,
une pondération de la value-at-risk en stress, un lookback long qui contient toujours une crise,
un coussin constitué en période calme et libéré en stress, et des limites sur la vitesse des
hausses. Tous relèvent la marge moyenne, que les membres paient chaque jour~; le dernier laisse
en outre la chambre de compensation sous-couverte le temps de rattraper son retard.
\iqlookfor{au moins deux outils et le coût de chacun.}
\end{solution}

\begin{solution}{iq:m1:margin:4}
Les scénarios extrêmes et le minimum sur les options vendues. Dans l'\omterm{def:m1:margin:scan}{intervalle de balayage}
normal, une option très hors de la monnaie ne perd presque rien~; un balayage seul laisserait
donc un vendeur accumuler une taille illimitée pour une marge presque nulle,
une position qui perd de façon catastrophique précisément dans le mouvement que le balayage
n'atteint pas. Les scénarios extrêmes réévaluent à deux ou trois intervalles et comptent
une fraction de la perte~; le minimum place un plancher sous chaque option vendue.
\iqlookfor{pourquoi les scénarios standard sont aveugles à la position.}
\end{solution}

\begin{solution}{iq:m1:margin:5}
Suivre les liquidités, et non seulement le P\&L~: la \omterm{def:m1:clearing-and-settlement:margin}{marge de variation} quotidienne, la \omterm{def:m1:clearing-and-settlement:margin}{marge initiale}
selon la méthode de la chambre de compensation avec ses paramètres tels qu'ils étaient à
chaque date (les taux de marge montent en période de crise, précisément lorsqu'une stratégie
à effet de levier perd), le multiplicateur du courtier, et les intérêts sur le collatéral déposé.
Imposer la contrainte que la trésorerie ne devienne jamais négative~; dimensionner les positions sur
la marge en stress, et non sur la marge courante. Rapporter le rendement sur le capital effectivement
mobilisé, coussin compris.
\iqlookfor{les taux de marge historiques et la contrainte de trésorerie.}
\end{solution}

\begin{solution}{iq:m1:margin:6}
Le P\&L se compense entre les deux maisons~; les liquidités, non. Lors d'un mouvement important, l'une des
maisons appelle la \omterm{def:m1:clearing-and-settlement:margin}{marge de variation} sur la jambe perdante dès le matin, tandis que le gain réalisé
auprès de l'autre est versé selon le cycle de cette dernière, éventuellement plus tard et dans une autre
devise. La \omterm{def:m1:clearing-and-settlement:margin}{marge initiale} augmente chez les deux. Dimensionner le coussin à partir d'un stress conjoint~:
le plus grand mouvement sur plusieurs jours du facteur de risque commun, chaque modèle de maison
réexécuté à la volatilité en stress, auquel s'ajoutent l'écart de temps entre le paiement d'un côté et
la réception de l'autre, ainsi que le risque de change. Pré-organiser les lignes de crédit et la transformation
de collatéral~; tester opérationnellement l'\omterm{def:m1:margin:intraday}{appel de marge intrajournalier}.
\iqlookfor{l'écart de temps et la hausse simultanée de la \omterm{def:m1:clearing-and-settlement:margin}{marge initiale}.}
\end{solution}
